\section{Systematization of Anti-Distillation Defenses}
\label{sec:similarity-attribution}

\begin{table*}[!t]
\centering
\caption{Comparison of representative defenses. \textcolor[HTML]{D46A6A}{Light red text} marks our findings under the conditions in Section~\ref{sec:systematization-findings}.}
\label{tab:anti-distillation-comparison}
\begin{threeparttable}
\scriptsize
\setlength{\tabcolsep}{2.5pt}
\renewcommand{\arraystretch}{1.12}
\begin{tabular}{@{}
>{\raggedright\arraybackslash}p{0.12\textwidth}
>{\raggedright\arraybackslash}p{0.15\textwidth}
>{\raggedright\arraybackslash}p{0.16\textwidth}
>{\raggedright\arraybackslash}p{0.11\textwidth}
>{\raggedright\arraybackslash}p{0.18\textwidth}
>{\raggedright\arraybackslash}p{0.22\textwidth}@{}}
\noalign{\hrule height 1.5pt}
\multirow{2}{*}{\textbf{Method}} & \multirow{2}{*}{\textbf{Signal}} &
\multicolumn{3}{c}{\textbf{Access and information requirements}} &
\multirow{2}{*}{\textbf{Main limitation}} \\
\cline{3-5}
& & \textbf{Provider side} & \textbf{Student side} & \textbf{Extra} & \\
\hline
\multicolumn{6}{l}{\textit{Traffic screening: evidence of collection}} \\
Provider defenses~\cite{OpenAIkey,anthropicDetectingPreventing} &
Request and account patterns &
Service logs & None & Account metadata &
Collection does not establish training. \\
\hline
\multicolumn{6}{l}{\textit{Active watermarking: provider-controlled signals}} \\
ADFP~\cite{xu2026antidistillation} &
Token preferences &
Control token sampling & Outputs & Proxy model; key; task queries &
Dilution; \textcolor[HTML]{D46A6A}{tool-call utility loss}. \\
AgentWM~\cite{wang2026protecting} &
Tool-variant frequencies &
Control tool selection & Outputs & Target distributions; domain tasks &
\textcolor[HTML]{D46A6A}{Limited transfer and dilution sensitivity}. \\
SeqWM~\cite{an2026sequential} &
Action sequences &
Control action selection & Outputs & Key; action vocabulary &
\textcolor[HTML]{D46A6A}{No tested student passes detection}. \\
\rowcolor{gray!12}
\name{} &
Context-conditioned behavior &
Insert tool calls; screen traffic & Outputs & Private trajectory records; scorer and safety models &
Requires traffic screening and retained evidence. \\
\hline
\multicolumn{6}{l}{\textit{Passive fingerprinting: existing model similarities}} \\
PoS Templates~\cite{wadhwa2025taught} &
Syntactic patterns &
Teacher outputs & Outputs & Candidate teachers; PoS templates &
Shared data and later training affect the signal. \\
Data-Free KD~\cite{shi2026knowledge} &
Internal model statistics &
Teacher API & Weights & Synthesized queries &
Requires student weights. \\
DistillDetect~\cite{rawat2026reference} &
Relative likelihoods &
Teacher outputs & Token logprobs & Reference models; calibration data &
Requires token scores and reference calibration. \\
SCOUT~\cite{jang2026retrospective} &
Output patterns &
Teacher outputs & Outputs & Candidate teachers; shared prompts &
Shared patterns; overwriting by later training. \\
AgentEcho~\cite{yang2026agents} &
Action-graph similarity &
Teacher trajectories & Outputs & Task harness; graph extraction &
\textcolor[HTML]{D46A6A}{Dependency-score collisions; nonspecific similarity gains}. \\
Reasoning PF~\cite{panfilov2026stealing} &
Answer overlap after reasoning prefill &
Teacher reasoning traces & Outputs; reasoning prefill & Reference answers &
\textcolor[HTML]{D46A6A}{Overlap from shared solutions and presentation}. \\
\noalign{\hrule height 1.5pt}
\end{tabular}
\begin{tablenotes}[flushleft]
\scriptsize
\item[] Provider side lists required access or control over the teacher service;
student side lists the audit interface. Outputs include visible actions.
Token logprobs require scoring supplied text. Extra lists auxiliary resources
and retained information; it is not a cost ranking.
\end{tablenotes}
\end{threeparttable}
\end{table*}

% Writing guidance: define the classification dimensions, interpret the table,
% then separate our new findings from limitations reported in prior work.
% Keep experimental configurations, numerical results, and counterexamples
% in Appendix~\ref{sec:systematization-details}.
\subsection{Classification Framework}
\label{sec:classification-framework}

We examine existing defenses from the provider's perspective: what evidence
does each method produce, and what does that evidence establish?
We group methods by their intervention point and source of evidence.
\emph{Traffic screening} identifies suspicious collection during model serving.
\emph{Active watermarking} embeds controlled signals into teacher outputs for
later verification. \emph{Passive fingerprinting} compares existing model
properties without embedding a signal.
Our scope includes both methods designed to detect distillation and output
watermarks whose use for this purpose requires the signal to survive training.

\partitle{Comparison dimensions}
We distinguish evidence of \emph{collection}, \emph{output provenance}, and
\emph{student training}: detecting a watermark in a collected output does not
by itself show that a student learned it.
We separate provider-side access or control, the student audit interface,
and extra resources such as secret keys, retained records, or reference models.
Output access does not imply access to token scores, weights, or a reasoning-prefill interface.
For deployment, we examine changes to serving, additional computation, and
effects on task utility. For robustness, we examine whether the signal
survives changes in the student model and processing of the training data.
These dimensions describe operational requirements; they do not assign a
common numerical score to methods evaluated under different protocols.

\subsection{Systematization Results}
\label{sec:systematization-results}



Table~\ref{tab:anti-distillation-comparison} separates the evidence each family
provides from the conditions needed to interpret it.
Traffic screening can identify suspicious collection and guide access controls,
but does not observe whether the collected outputs enter student training.
Active watermarking supplies a provider-controlled signal. Its value for
distillation detection depends on whether the student learns that signal and
whether independent models reproduce it without exposure.
Passive fingerprinting avoids modifying teacher outputs, but shared patterns
require controls for other sources of similarity. Some methods also require
student logits or weights that a deployed API does not expose.

The comparison reveals a gap between detecting collection and attributing
student training. Neither an output watermark nor a high similarity score
alone resolves this gap. A practical defense must connect its evidence to
training while preserving utility and surviving data processing.
The surveyed results do not establish these properties jointly under our
serving and post-hoc audit setting.

\subsection{New Findings}
\label{sec:systematization-findings}

\partitle{Finding 1: Behavioral watermarks can lose detectability through
distillation and dilution}
We evaluate AgentWM and SeqWM with multiple student models on BFCL and
SWE-bench. AgentWM biases the choice among functionally equivalent tools;
SeqWM embeds patterns in action sequences.
Under our evaluation protocol, AgentWM transfers detectably to only a subset
of students, while no SeqWM student meets the model-level detection rule.
For AgentWM, both initially detectable settings lose detection after mixing
watermarked trajectories with an equal amount of clean data calibrated to
the unwatermarked tool distribution.
Thus, preserving useful task behavior during training need not preserve a
detectable watermark. These results concern the tested implementations and
decision rules; they do not establish that every watermark signal disappears.
Appendix~\ref{sec:watermark-limitations} gives the configurations, per-model
results, and a separate ADFP case study on tool-call utility.

\partitle{Finding 2: Similarity scores can give ambiguous evidence of a training source}
Passive fingerprinting uses similarities between models to infer a training
relationship. But the selected patterns may capture shared conventions or
miss meaningful differences. On two public reasoning problems, prefix changes
and conventional solutions yield substantial answer overlap under our local
n-gram measure. An alternative prefix can yield higher overlap than the
teacher prefix, while a standard mathematical solution can yield high overlap
even without a reasoning prefix.

The problem also extends to tool-use structure. For AgentEcho, we construct
graphs with very different dependencies that receive identical
dependency-similarity scores. Beyond these constructed examples, our
trajectory audit finds high scores across structurally different executions.

We also find that a reasoning-distilled model becomes more similar to several
comparison models. The largest increase is toward Kimi-K2, although the
model card names Claude Opus~4.6 as the teacher. Such changes can be difficult
to interpret against the overlapping training histories of modern models.
Appendices~\ref{sec:textual-similarity} and~\ref{sec:dependency-ambiguity}
provide the counterexamples, controls, and measurements.

These findings motivate a defense that retains private evidence during
serving and tests its inheritance after distillation. \name{} uses traffic
screening to select candidate interactions, then compares the suspect
student's behavior under recorded contexts and modified control contexts.
